\documentclass[11pt]{article}
\usepackage[margin=1in]{geometry}
\usepackage{amsmath,amssymb}
\setlength{\parindent}{0pt}
\setlength{\parskip}{6pt}
\begin{document}
\section*{STAT 412 QZ08 --- Question 6}

\subsection*{Solution}

\[
H_0:\text{whether a subject lies is independent of the polygraph indication.}
\]
\[
H_1:\text{whether a subject lies is not independent of the polygraph indication.}
\]

Expected counts:
\[
E=\frac{(\text{row total})(\text{column total})}{\text{grand total}}.
\]
For example, for ``indicated lie'' and ``did not lie,''
\[
E=\frac{64(22)}{80}=17.6.
\]
\[
\begin{array}{c|cc}
 & \text{Did not lie} & \text{Lied}\\ \hline
\text{Indicated lie} & 17.6 & 46.4\\
\text{Indicated no lie} & 4.4 & 11.6
\end{array}
\]

Then
\[
\chi^2=\sum\frac{(O-E)^2}{E}
=\frac{(15-17.6)^2}{17.6}
+\frac{(49-46.4)^2}{46.4}
+\frac{(7-4.4)^2}{4.4}
+\frac{(9-11.6)^2}{11.6}
=2.649.
\]
For a \(2\times2\) table,
\[
df=(2-1)(2-1)=1.
\]

\[
\chi^2=2.649,\qquad df=1,\qquad p=0.1036.
\]

Since \(p>0.05\), fail to reject \(H_0\). The results do not suggest that polygraphs are effective in distinguishing between truth and lies.

\subsection*{Final Answer}

\fbox{\begin{minipage}{0.94\linewidth}
Expected table: \(17.6,46.4; 4.4,11.6\). \(\chi^2=2.649\), \(p=0.1036\). Fail to reject \(H_0\); no, the results do not suggest effectiveness.
\end{minipage}}

\end{document}
